\documentclass{situationist-preprint}

\title{The Quus Hypothesis Refuted}
\msid{LNG-05-b}
\msfamily{LNG-05}
\msdiscipline{Philosophy of Language}
\mstype{Companion manuscript}
\mskeywords{rule-following; Kripke; quus; meaning scepticism; primitive meaning facts}

\begin{document}
\maketitle

\begin{abstract}
Kripke's sceptic proposes that by \enquote{plus} a speaker may in the past have meant quus, a function that agrees with addition on all arguments she has computed and diverges elsewhere. LNG-05-a argues that there is a primitive fact that the speaker means addition. This manuscript draws the consequence for the sceptic's specific hypothesis. Given that fact, the quus hypothesis is false, and it is shown to be false without any appeal to dispositions or to a reduction of meaning. We add that the hypothesis cannot be stated by a sceptic who accepts the reduction premise his argument needs, since that premise would deprive his own \enquote{quus} of a meaning.
\end{abstract}

\section{The hypothesis and its refutation}

Kripke defines quus as the function that takes $x$ and $y$ to $x+y$ if both are less than 57, and to 5 otherwise, and imagines a speaker who has never computed with numbers as large as 57 \citep{Kripke1982}. The sceptic's hypothesis is that by \enquote{plus} she meant quus all along; when she now answers \enquote{125} to \enquote{68 plus 57}, she departs from her past meaning.

Quus fits the speaker's base in the sense of Definition~2.2 of LNG-05-a, and so Assumption~2.4 (Fit) does not exclude it. That is the point of the example, and LNG-05-a grants it in Lemma~3.1. The refutation must come from elsewhere.

\begin{proposition}\label{p:refute}
Given Assumptions~2.5 (Determinacy) and~2.7 (Datum) of LNG-05-a, the speaker of the datum does not mean quus by \enquote{plus}.
\end{proposition}

\begin{proof}
By the datum, she means addition. Addition and quus are distinct functions, since $68 + 57 = 125$ and $68 \oplus 57 = 5$. By Determinacy she means at most one function. So she does not mean quus.
\end{proof}

The proof is short because the work is done by the datum, which LNG-05-a defends as a primitive fact (Theorem~3.3 there shows that Reduction is false given it). The refutation is of the same form as the refutation of any sceptical hypothesis by a known fact incompatible with it, and its strength is the strength of our knowledge that we mean addition. What it does not need is a fact about the speaker, described in non-semantic terms, that singles out addition; LNG-05-a argues that there is none, and that none is required.

\section{The hypothesis undermines itself}

The sceptic grants, for the sake of the argument, that his present words have their meanings; he asks only whether the speaker's past words did \citep{Kripke1982}. His argument against past meaning, however, depends on Reduction (Definition~2.6 of LNG-05-a): only a base that singles out a function could constitute a meaning fact. Reduction is general. Applied to the sceptic's present use of \enquote{quus}, it yields, by Theorem~3.2 of LNG-05-a, that he means no function by \enquote{quus}, since his base for that word is also finite and fitted by deviant functions. The hypothesis he states then has no determinate content, and there is no determinate alternative to addition for the speaker to have meant. The point is a special case of Boghossian's observation that the sceptical argument, if sound, generalizes to all content, including the content of the sceptic's own thoughts \citep{Boghossian1989}.

So the sceptic faces a dilemma. If he grants himself determinate present meaning while denying it to the speaker's past, his premise is applied selectively. If he applies it uniformly, his own hypothesis has no content. On the first horn the argument is unsound; on the second, there is no quus hypothesis to refute.

\begin{objection}
Proposition~\ref{p:refute} begs the question, since the datum is precisely what the sceptic denies.
\end{objection}

\begin{reply}
The sceptic does not offer independent grounds against the datum; his ground is the argument from Reduction, whose premise LNG-05-a shows to be false given the datum and unsupported without it (Proposition~3.4 there). A refutation that relies on a premise the opponent rejects begs the question only if the premise has no support beyond the conclusion. The datum is supported by everything we know about our own arithmetical practice, and the rival argument is unsound.
\end{reply}

\section{Conclusion}

The quus hypothesis is consistent with the speaker's base, and false. It is refuted by the primitive fact that she means addition, together with the determinacy of meaning, and a sceptic who accepted the reduction his argument requires could not state it.

\begin{thebibliography}{9}\small
\bibitem[Boghossian(1989)]{Boghossian1989} Boghossian, P.~A. (1989). The rule-following considerations. \emph{Mind}, 98, 507--549.
\bibitem[Kripke(1982)]{Kripke1982} Kripke, S.~A. (1982). \emph{Wittgenstein on Rules and Private Language}. Harvard University Press.
\end{thebibliography}

\end{document}
